%!TEX root = ../main.tex
\section{Introduction}
\label{sec:introduction}

Robot learning is entering a data-rich but not yet data-coherent regime.
Open X-Embodiment aggregates heterogeneous datasets into a common interface.
DROID distributes real-world collection across institutions, while the Universal
Manipulation Interface (UMI) moves
demonstration away from a continuously occupied robot.  Recent generalist
policies train on mixtures that span embodiments, tasks, and supervision
types~\cite{collaboration2023open_2310_08864,khazatsky2024droid_2403_12945,chi2024universal_2402_10329,kim2024openvla_2406_09246,team2024octo_2405_12213}.
This expansion is often summarized by episode counts or hours.  Neither is a
sufficient statistic.  A million records can contain incompatible coordinate
frames, control modes, frequencies, language granularity, success criteria,
operator styles, and physical assumptions.  Even internally valid records can be
redundant for one learner, harmful to another, or useful only at a particular
training stage.

The missing object of study is the transformation between a raw interaction and
the distribution actually sampled by an optimizer.  In this survey, \emph{data
preparation for embodied intelligence} denotes the methods that design or import
data sources and subsequently validate, repair, align, segment, annotate, score,
filter, weight, route, augment, synthesize, verify, package, and iteratively evolve
their records.  Its output is not merely a cleaned dataset.  It is a versioned
training distribution conditioned on a target learner, embodiment, task
distribution, training stage, deployment environment, and budget.
When feedback from a trained or deployed model changes a later data decision,
we call the resulting closed loop a \emph{data engine}.  This definition is
narrower than using the term for any data-processing pipeline.
When such feedback repeatedly revises the training distribution across rounds,
we refer to the process as \emph{data self-evolution}.

This boundary is deliberately narrower than ``data for embodied AI.''  Our core
scope is manipulation and generalist vision--language--action (VLA) learning.  We
include collection when the acquisition design changes compatibility, coverage,
or feedback signals.  Human video and simulation enter the review when they are
transformed into robot-learning supervision.  Model and system reports enter
when they disclose a data operation or recipe.  We do not systematically review
robot hardware, policy architectures, generic simulators, navigation, autonomous
driving, or benchmarks
that only evaluate a final model.  These topics enter only when they expose a
general data-preparation mechanism.

\subsection{Why a Separate Survey Is Needed}

Adjacent reviews provide valuable maps, but their primary units differ from ours
(Table~\ref{tab:adjacent-surveys}).  The Data Pyramid is source-centric.  A recent
VLA survey catalogs datasets, benchmarks, and data engines, while the human-video
survey organizes transfer mechanisms~\cite{ye2026data_2607_24744,wang2026vision_2604_23001,ma2026robot_2604_27621}.
They do not center the decision problem common to all sources: \emph{which
transformation should be applied, using which signal, and what evidence shows
that it increased training value?}  That decision is becoming a distinct
technical layer for three reasons.

Data operations increasingly contain learned components.  Recent methods learn
source-mixture weights, estimate redundancy and information, predict task
progress, or approximate how a training record changes a target learner
~\cite{hejna2024re_2408_14037,hejna2025robot_2502_08623,dass2025datamil_2505_09603,zhang2025scizor_2505_22626,xu2026athena_2606_16208}.
Foundation models can also label or generate data at scale, which moves the
bottleneck from production to calibration and verification
~\cite{blank2024scaling_2410_17772,mandlekar2023mimicgen_2310_17596,li2026xiaomi_2607_11643}.
Meanwhile, self-improving systems turn deployment outcomes into future collection and
post-training data, so preparation becomes a feedback policy rather than a
one-shot preprocessing script~\cite{bousmalis2023robocat_2306_11706,ahn2024autort_2401_12963,intelligence20250_2511_14759,pan2026sop_2601_03044,wang2026learning_2605_00416}.
The same shift is visible in industrial reports.  Qwen foregrounds
representation, motion, and behavior alignment.  The $\pi$ family describes
stage- and metadata-dependent routing.  GR00T combines real anchors with generated
and simulated expansion.  Gemini reports Motion Transfer, a training mechanism
that aligns motion across robot embodiments
~\cite{yuan2026qwen_2606_17846,intelligence20260_2604_15483,nvidia2025gr00t_2503_14734,team2025gemini_2510_03342}.
These reports do not support a leaderboard, but they show that data preparation
has become a recurring systems concern rather than a hidden preprocessing step.

\begin{table*}[t]
\centering
\caption{Difference in organizing focus from adjacent surveys. A check mark indicates a central organizing dimension rather than incidental coverage.}
\label{tab:adjacent-surveys}
\small
\setlength{\tabcolsep}{5pt}
\begin{tabularx}{\textwidth}{@{}>{\RaggedRight\arraybackslash}X>{\RaggedRight\arraybackslash}p{3.4cm}cccc@{}}
\toprule
Survey & Organizes the field by & Curation & Utility & Feedback & Protocol \\
\midrule
Data Pyramid~\cite{ye2026data_2607_24744}
& Data-source hierarchy & \cmark & \xmark & \xmark & \xmark \\
VLA data survey~\cite{wang2026vision_2604_23001}
& Dataset and engine type & \cmark & \xmark & \cmark & \xmark \\
Human-video survey~\cite{ma2026robot_2604_27621}
& Transfer path & \xmark & \xmark & \xmark & \xmark \\
\textbf{This survey}
& Preparation decision & \cmark & \cmark & \cmark & \cmark \\
\bottomrule
\end{tabularx}
\end{table*}


\subsection{Central Thesis}

Our central thesis is that embodied data quality is a \emph{conditional,
multidimensional utility}, not an intrinsic scalar.  Consider a failed trajectory.
It may be harmful if its action labels are silently treated as expert targets,
valuable if its successful prefix or recovery is segmented, and essential if it
is the only observation of a deployment failure.  Conversely, an apparently
smooth successful trajectory can add little if it duplicates states already
covered or encodes an action convention incompatible with the target robot.
Formal analyses of imitation data already show that state diversity, action
divergence, and uncertainty in state transitions interact rather than admit a universal
ordering~\cite{belkhale2023data_2306_02437}.  Empirical curation work reinforces
the point: removing half of a dataset can help in one setting but lose important
coverage in another~\cite{hejna2025robot_2502_08623,bedi2026what_2606_10229}.

This perspective leads to four orthogonal axes.  The \emph{operation axis} asks
what changes in the record or distribution.  The \emph{signal axis} asks what
evidence guides that change.  The \emph{utility-evidence axis} asks whether the
claim is supported by an intrinsic diagnostic, a proxy, controlled policy
training, or deployment.  The \emph{lifecycle axis} asks whether a decision is
static, training-time, post-training, or continually revised.  Source type,
modality, transformed object, granularity, and objective are recorded as additional attributes,
not used as mutually exclusive bins.

\subsection{Contributions and Reader's Guide}

The survey makes four contributions.  First, it defines data preparation as the
transformation of raw interactions into target-conditional training
distributions, separating this problem from general data collection and policy
architecture design.  Second, it provides an evidence map of 211 works with
search records, reading-depth labels, and structured coding.  The resulting
taxonomy compares operations by their decision signal, granularity, cost,
validation, generalization, and failure modes instead of headline performance.
Third, the survey reconstructs public industrial recipes and compares data
engines by the feedback they observe, the data decision they make, and the
controls used during model updates.  This comparison also separates disclosed
facts from within-paper evidence and cross-paper interpretation.  Finally, the
survey relates embodied preparation to language-model data engineering and
develops a research agenda for utility calibration, stable closed-loop updates,
and bounded preparation autonomy.

Section~\ref{sec:method} describes the review protocol.  Section~\ref{sec:problem}
formalizes the problem and taxonomy.  Sections~\ref{sec:acquisition}--
\ref{sec:evolution} review the technical lifecycle.
Section~\ref{sec:recipes-engines} synthesizes industrial recipes, post-training
loops, and the autonomy ladder.  Section~\ref{sec:evaluation} examines evidence,
lessons from large language model data preparation, and cross-cutting
conclusions.  Section~\ref{sec:agenda} then develops the research agenda.
